\input{../common}
\input{../bibliography}

\begin{document}
  %<*content>
  \development{algebra}{theoreme-de-wedderburn}{Théorème de Wedderburn}

  \summary{En utilisant les polynômes cyclotomiques, nous montrons que tout corps fini est commutatif.}

  \reference[GOU21]{100}

  \begin{lemma}
    \label{theoreme-de-wedderburn-1}
    Soient $\mathbb{K}$ et $\mathbb{L}$ deux corps finis tels que $\mathbb{K}$ est commutatif et $\mathbb{K} \subseteq \mathbb{L}$. Alors $\exists d \in \mathbb{N}^*$ tel que $|\mathbb{L}| = |\mathbb{K}|^d$.
  \end{lemma}

  \begin{proof}
    $\mathbb{L}$ est un espace vectoriel sur $\mathbb{K}$ de dimension finie $d$ (car $\mathbb{L}$ est fini). Donc $\mathbb{L}$ est isomorphe en tant que $\mathbb{K}$-espace vectoriel à $\mathbb{K}^d$. En particulier, $|\mathbb{L}| = |\mathbb{K}|^d$.
  \end{proof}

  \begin{theorem}[Théorème de Wedderburn]
    Tout corps fini est commutatif.
  \end{theorem}

  \begin{proof}
    Soit $\mathbb{K}$ un corps. L'idée va être de procéder par récurrence sur le cardinal du corps.
    \begin{itemize}
      \item \uline{Si $|\mathbb{K}| = 2$ :} alors $\mathbb{K} = \{0, 1\}$ est commutatif.
      \item \uline{On suppose le résultat vrai pour tout corps fini de cardinal strictement inférieur à $\mathbb{K}$.} On veut montrer que $\mathbb{K}$ est commutatif. Supposons par l'absurde que $\mathbb{K}$ ne l'est pas. On pose
      \[ Z = Z(\mathbb{K}) = \{ x \in \mathbb{K} \mid \forall y \in \mathbb{K}, \, xy = yx \} \]
      le centre de $\mathbb{K}$ dont on note $q$ le cardinal. C'est un sous-corps de $\mathbb{K}$ qui est (par hypothèse) inclus strictement dans $\mathbb{K}$. Donc $Z$ est commutatif, et par le \cref{theoreme-de-wedderburn-1}, on peut écrire $|\mathbb{K}| = q^n$ où $n \in \mathbb{N}^*$. Si $x \in \mathbb{K}$, on pose
      \[ \mathbb{K}_x = Z_{\mathbb{K}}(\{ x \}) = \{ y \in \mathbb{K} \mid xy = yx \} \]
      Montrons que
      \[ \exists d \mid n \text{ tel que } |\mathbb{K}_x| = q^d \tag{$*$} \]
      Notons déjà encore une fois que $\mathbb{K}_x$ est un sous-corps de $\mathbb{K}$.
      \begin{itemize}
        \item Si $\mathbb{K}_x = \mathbb{K}$, on a $|\mathbb{K}_x| = |\mathbb{K}| = q^n$. Il suffit donc de prendre $d = n$.
        \item Sinon, $\mathbb{K}_x \subsetneq \mathbb{K}$, donc $\mathbb{K}_x$ est commutatif par hypothèse. Par le \cref{theoreme-de-wedderburn-1}, il existe $k \in \mathbb{N}^*$ tel que $|\mathbb{K}| = |\mathbb{K}_x|^k$.
        \newpar
        Mais, $Z$ est un sous-corps (commutatif) de $\mathbb{K}_x$, donc d'après le \cref{theoreme-de-wedderburn-1}, il existe $d \in \mathbb{N}^*$ tel que $|\mathbb{K}_x| = |Z|^d$. Donc on a
        \[ q^n = |\mathbb{K}| = |\mathbb{K}_x|^k = (q^d)^k = q^{dk} \]
        d'où $d \mid n$.
      \end{itemize}
      On considère l'action par conjugaison de $\mathbb{K}^*$ sur lui-même $(x, y) \mapsto xyx^{-1}$. Si $y \in \mathbb{K}^*$, alors
      \[ \stab_{\mathbb{K}^*}(y) = \{ x \in \mathbb{K}^* \mid x.y = y \} = \mathbb{K}_y^* \]
      Soit $\Omega$ un système de représentants associé à la relation d'équivalence ``être dans la même orbite''. L'équation aux classes donne alors
      \[ |\mathbb{K}^*| = \sum_{\omega \in \Omega} \frac{|\mathbb{K}^*|}{|\stab_{\mathbb{K}^*}(\omega)|} \]
      Or,
      \[ \stab_{\mathbb{K}^*}(\omega) = \mathbb{K}^* \iff \forall x \in \mathbb{K}^*, \, \omega x = x \omega \iff \omega \in Z^* \]
      donc en notant $\Omega' = \Omega \setminus Z^*$, on a :
      \[ |\mathbb{K}^*| = \sum_{\omega \in Z^*} \frac{|\mathbb{K}^*|}{|\stab_{\mathbb{K}^*}(\omega)|} + \sum_{\omega \in \Omega'} \frac{|\mathbb{K}^*|}{|\stab_{\mathbb{K}^*}(\omega)|} = |Z^*| + \sum_{\omega \in \Omega} \frac{|\mathbb{K}^*|}{|\mathbb{K}^*_\omega|} \tag{$**$} \]
      Soit $\omega \in \Omega'$. Par $(*)$,
      \[ \exists d \mid n \text{ tel que } |\stab_{\mathbb{K}^*}(\omega)| = |\mathbb{K}^*_\omega| = q^d - 1 \]
      De plus, $d \neq n$ (car $\omega \notin Z^*$). Si maintenant on pose
      \[ \forall d \mid n, \, \lambda_d = |\{ \omega \in \Omega' \mid |\stab_{\mathbb{K}^*}(\omega)| = q^d - 1 \}| \]
      on peut alors écrire en remplaçant dans $(**)$ :
      \[ q^n - 1 = |K^*| = (q - 1) + \sum_{d \mid \mid n} \lambda_d \left( \frac{q^n - 1}{q^d - 1} \right) \tag{$***$} \]
      Si $d \mid \mid n$, on a
      \[ X^n-1 = \prod_{k \mid n} \Phi_k = \Phi_n \left ( \prod_{k \mid d} \Phi_k \right ) \left ( \prod_{\substack{k \mid \mid n \\ k \nmid d}} \Phi_k \right ) = \Phi_n (X^d - 1) \left ( \prod_{\substack{k \mid \mid n \\ k \nmid d}} \Phi_k \right ) \]
      Donc, $\Phi_n \mid \frac{X^n - 1}{X^d - 1}$ dans $\mathbb{Z}[X]$. Ceci étant vrai quel que soit $d$ divisant strictement $n$, on en déduit
      \[ \Phi_n \mid \sum_{d \mid \mid n} \lambda_d \frac{X^n - 1}{X^d - 1} \text{ dans } \mathbb{Z}[X] \]
      Comme de plus, $\Phi_n \mid X^n - 1$ dans $\mathbb{Z}[X]$, on conclut que
      \[ \Phi_n \mid X^n - 1 - \sum_{d \mid \mid n} \lambda_d \frac{X^n - 1}{X^d - 1} \text{ dans } \mathbb{Z}[X] \]
      ce qui donne, une fois évalué en $q$ :
      \[ \Phi_n(q) \mid q^n - 1 - \sum_{d \mid \mid n} \lambda_d \frac{q^n - 1}{q^d - 1} \overset{(***)}{=} q-1 \implies |\Phi_n(q)| \leq q-1 \]
      Mais $n \geq 2$, donc
      \begin{align*}
        |\Phi_n(q)| &= \prod_{\xi \in \mu_n^*} |q - \xi| \\
        &> \prod_{i=1}^{\varphi(n)} |q - 1| \\
        &\geq |q-1|
      \end{align*}
      On peut en effet interpréter $|q - \xi|$ comme la distance du complexe $q$ au complexe $\xi$ ; le premier est sur l'axe réel et est supérieur ou égal à $2$, le second est sur le cercle unité mais n'est pas sur l'axe réel :
      \begin{center}
        \begin{tikzpicture}
          \draw[->] (-3, 0) -- (5, 0) node[right] {$x$};
          \draw[->] (0, -3) -- (0, 3.5) node[above] {$y$};
          \draw[orange, thick] (2, 0) -- (4, 0) node[below, shift={(-1,0)}]{$|1-q|$};
          \draw[teal, thick] (45:2) -- (4, 0) node[above, shift={(-1.2,0.6)}, rotate=-27.5]{$|\xi-q|$};
          \draw(0,2) node {$\bullet$};
          \draw(0,2) node[above right]{$i$};
          \draw(2,0) node {$\bullet$};
          \draw(2,0) node[below right]{$1$};
          \draw(0,0) circle (2);
          \draw(4,0) node {$\bullet$} node[below]{$q$};
          \draw(45:2) node {$\bullet$} node[above right]{$\xi$};
        \end{tikzpicture}
      \end{center}
      cela nous permet de justifier l'inégalité stricte. On a donc une contradiction.
    \end{itemize}
  \end{proof}
  %</content>
\end{document}
